



















Min LIA har gett mig värdefull erfarenhet av både tekniska verktyg och arbetsprocesser. Jag har fått förståelse för hur ett team samarbetar i verkliga projekt, hur man hanterar tickets, och hur man följer agil metodik. 

Jag känner att jag har vuxit i min förmåga att felsöka problem, planera arbetsuppgifter och kommunicera med kollegor. LIA-perioden har stärkt mitt intresse för DevOps och molnlösningar.


\section{Reflektion över kunskaper och färdigheter}
\label{sec:reflektion}

Under min LIA-period har jag fått mycket stöd och hjälp från både min handledare och resten av teamet på Insighta. De har gett mig tydliga uppgifter, tid för att lära mig, samt tillgång till de plattformar och verktyg jag behövde, till exempel GCP och GKE. Jag fick även en ny MacBook för att kunna arbeta effektivt. Jag har blivit bemött som en riktig medarbetare och det har hjälpt mig att växa både tekniskt och i min arbetsroll.

Vår kommunikation skedde främst via Slack och e-post, och vi hade regelbundna möten. Varje vecka hade teamet också ett tekniskt möte där vi lärde oss om nya verktyg, arbetssätt och lösningar som kan användas inom marketing analytics. Detta gav mig mycket inspiration och ökade min förståelse för hur ett professionellt DevOps-arbete går till.


\subsection{Feedback som jag har fått gällande mina kunskaper och färdigheter}
\label{subsec:feedback}
\subsubsection{Kunskaper}
\label{subsubsec:kunskaper_feedback}

Jag har fått positiv feedback på min vilja att lära mig nya tekniska verktyg och på att jag snabbt utvecklade min förståelse för Kubernetes, GCP och Terraform/CDKTF, trots att dessa områden inte ingick i våra kurser på skolan. Teamet uppskattade att jag ställde frågor när jag behövde det och att jag tog ansvar för att förstå helheten i systemen.

\subsubsection{Omdömes-, planerings- och prioriteringsförmåga}
\label{subsubsec:omdome_feedback}

Min handledare sa att jag var bra på att planera mina arbetsuppgifter. Jag fick också höra att jag tog mig an utmanande problem på ett strukturerat sätt, särskilt när jag felsökte Dagster och Kubernetes. Jag har också fått beröm för att jag håller fokus, även när prioriteringar ändras.

\subsubsection{Samarbets- och kommunikationsförmåga}
\label{subsubsec:samarbete_feedback}

Jag har fått positiv feedback kring hur jag kommunicerar och samarbetar. Jag har varit aktiv i Slack-kanalerna, haft bra dialog med teamet och deltagit i möten. Jag visade att jag kan ta emot och ge feedback, och att jag kan samarbeta även på distans, vilket är viktigt eftersom teamet är internationellt.

\subsubsection{Prestationsförmåga och målfokusering}
\label{subsubsec:prestation_feedback}
Jag har fått beröm för att jag alltid försöker göra mitt bästa och att jag är målmedveten. När jag fick min första svårare ticket med Dagster och Kubernetes jobbade jag metodiskt tills jag uppfyllde alla acceptance criteria. De uppskattade också att jag tog initiativ och växte in i rollen steg för steg.


\subsection{Hur jag har utvecklat mina färdigheter och kunskaper under terminen}
\label{subsec:utveckling}
\subsubsection{Kunskaper}
\label{subsubsec:kunskaper_utveckling}
Jag har lärt mig mycket mer än vad jag förväntade mig. Jag har utvecklat tekniska kunskaper inom Kubernetes, GKE, Terraform/CDKTF, GCP, Helm, Dagster och CI/CD. Jag har också blivit mer självsäker i att felsöka problem och förstå stora system. Detta är kunskaper som inte täcktes i våra kurser, så allt detta har jag lärt mig direkt genom praktiskt arbete hos Insighta.

\subsubsection{Omdömes-, planerings- och prioriteringsförmåga}
\label{subsubsec:omdome_utveckling}
Jag har blivit bättre på att planera mitt arbete och bryta ned större uppgifter i mindre delar. Jag har lärt mig att prioritera det som är viktigast, och jag har blivit bättre på att hantera situationer där saker inte fungerar direkt. Detta märktes särskilt när jag arbetade med Dagster-deployen och när jag satte upp GKE-klustret.


\subsubsection{Samarbets- och kommunikationsförmåga}
\label{subsubsec:samarbete_utveckling}

Jag utvecklades mycket i kommunikationen eftersom jag arbetade i ett internationellt team. Jag blev bättre på att förklara mina problem, fråga om hjälp, och berätta hur långt jag kommit i mina uppgifter. Jag lärde mig också hur viktigt det är att skriva tydliga meddelanden, delta i möten och ge uppdateringar i tid.


\subsubsection{Prestationsförmåga och målfokusering}
\label{subsubsec:prestation_utveckling}

Under LIA-perioden blev jag mycket bättre på att hålla fokus på målen och jobba steg för steg tills uppgifterna var klara. Jag lärde mig att inte ge upp när saker blev svåra, utan istället testa olika lösningar tills det fungerade. Jag utvecklade en starkare förståelse för mitt eget ansvar och hur mitt arbete påverkar teamet.


\subsection{Hur jag ska förbättra mina kunskaper och färdigheter}
\label{subsec:forbattring}
Jag vill fortsätta fördjupa mig i Terraform, Kubernetes, GCP och CI/CD eftersom detta är centrala områden för DevOps. Jag vill också förbättra min förståelse för Dagster och börja gå mer mot MLOps, eftersom Insighta arbetar mycket med data-pipelines. Jag planerar att fortsätta öva på dessa verktyg under min andra LIA och även studera dem vid sidan av skolan, till exempel genom kurser, dokumentation och egna projekt.